\documentclass{article}
\begin{document}
\section*{Abstract}
\noindent Declining competition volume foreshadows exit from organised youth sport; whether the athlete's own performance development adds information is untested in population data. Mastery-experience accounts of motivation predict that visible personal progress retains athletes, most strongly where rank-based competence is scarce, as for athletes born late in the selection year. Using the Norwegian Athletics Federation's registry, we followed athletes born 1998--2005 with scorable age-graded points in their main event group at ages 13--14 and 15--16 (n = 3\,877; 53\,\% female), of whom 654 (16.9\,\%) were later active as seniors (age $\geq$20). Progression was the change in age-graded points (mean of three best scores) between the two windows. In logistic regression adjusted for sex, cohort, performance level, competition volume and the number of scored results, progression predicted senior activity (odds ratio [OR] 2.24 per SD, 95\,\% CI 1.90--2.63) and modestly improved discrimination (cross-validated AUC 0.854 to 0.866). The association replicated across birth cohorts and held at equal numbers of results, with event-equated points, and among athletes still competing after age 16 (OR 1.83). Aligned at exit, dropouts held their performance level until the final season and fell about 12 points below it as their competition volume collapsed. The pre-specified progression $\times$ birth-quarter interaction was null (OR 0.95, 95\,\% CI 0.84--1.06), and evidence for a mediated pathway was weak. Low progression accompanies and precedes exit, beyond volume and level, but the relative-age-specific form of the mastery prediction was not supported. Current form, not early rank, is what the retention signal tracks.

\vspace{0.5em}
\noindent\textbf{Keywords:} youth sport; dropout; athlete retention; performance development; mastery experience; relative age; track and field
\end{document}
